\documentclass[12pt, a4paper]{article}

% 引入必要的宏包
\usepackage[UTF8]{ctex} % 支持中文
\usepackage{geometry} % 页面边距设置
\geometry{left=2.5cm, right=2.5cm, top=3cm, bottom=3cm}
\usepackage{graphicx} % 插入图片
\usepackage{float} % 图片位置控制
\usepackage{caption} % 图片标题
\usepackage{subcaption} % 子图
\usepackage{color} % 颜色
\usepackage{booktabs} % 三线表
\usepackage{hyperref} % 超链接
\usepackage{fancyhdr} % 页眉页脚
\usepackage{amsmath}
% 设置超链接样式
\hypersetup{
    colorlinks=true,
    linkcolor=blue,
    filecolor=magenta,      
    urlcolor=cyan,
}

% 页眉页脚设置
\pagestyle{fancy}
\fancyhf{}
\fancyhead[L]{进阶流水灯控制系统设计与仿真}
\fancyhead[R]{实验考核报告}
\fancyfoot[C]{\thepage}

\usepackage{listings}
\usepackage{xcolor}

\definecolor{codegreen}{rgb}{0,0.6,0}
\definecolor{codegray}{rgb}{0.5,0.5,0.5}
\definecolor{codepurple}{rgb}{0.58,0,0.82}
\definecolor{backcolour}{rgb}{0.95,0.95,0.92}

\lstset{
    backgroundcolor=\color{backcolour},   
    commentstyle=\color{codegreen}\itshape,
    keywordstyle=\color{blue}\bfseries,
    numberstyle=\tiny\color{codegray},
    stringstyle=\color{codepurple},
    basicstyle=\ttfamily\footnotesize,
    breakatwhitespace=false,         
    breaklines=true,                 
    captionpos=b,                    
    keepspaces=true,                 
    numbers=left,                    
    numbersep=5pt,                  
    showspaces=false,                
    showstringspaces=false,
    showtabs=false,                  
    tabsize=2,
    frame=single,
    frameround=tttt,
    rulecolor=\color{black!30}
}
\setlength{\headheight}{15pt}
\begin{document}

% ---------------- 封面部分 ----------------
\begin{titlepage}
    \begin{center}
        \vspace*{2cm}
        \Huge{\textbf{智能自适应流水灯控制系统设计（PCB设计部分}}\\
        \vspace{0.5cm}
        \Large{\textbf{——基于555定时器与CD4017的综合应用与拓展，引入基于 Atmel AVR 架构的微控制器}}\\
        \vspace{2cm}
        
        \large
        \begin{tabular}{rl}
            \textbf{课程名称：} & 电子电路设计与仿真 \\
            \textbf{学生姓名：} & [姓名] \\  % <--- 这里改成你的名字
            \textbf{学号：} & [学号] \\      % <--- 这里改成你的学号
            \textbf{完成日期：} & \today \\
        \end{tabular}
        \vspace{3cm}
        
        \rule{0.8\textwidth}{0.5pt} \\
        \vspace{0.5cm}
        \textit{本项目以CD4017与555定时器为基础构建流水灯硬件架构。为提升系统的工程实用价值，分别通过引入模拟传感器实现\textbf{环境联动报警}，以及引入微控制器（MCU）实现\textbf{PWM呼吸算法}与\textbf{UART串口通信}，完成了具有较高复杂度的物联网控制系统原型设计。}
    \end{center}
\end{titlepage}

\newpage
\newpage

% ---------------- 正文部分 ----------------
\section{选题背景与设计目标}
本课程设计围绕专业课中信号与系统以及数字信号处理相关理论展开。在专业课程中，图像通常被视为二维离散信号。为了更直观地理解图像的时频域变换、空间边缘提取以及区域特征分析，我设计并开发了这款基于 MATLAB App Designer 的交互式图像分析系统。

本设计的目标是实现一个功能完整的图像处理软件，具备以下核心能力：
\begin{enumerate}
    \item 支持图像的导入、自适应增强、伪彩色映射及无损导出。
    \item 结合频域分析理论，实现二维快速傅里叶变换的可视化。
    \item 利用事件驱动机制，实现鼠标拖拽区域的实时截图与边缘轮廓同步提取。
    \item 重构图像空间灰度分布，生成可3D旋转的亮度拓扑图。
\end{enumerate}

\section{系统总体设计与界面布局}
为保障软件的易用性，我采用了单窗多选项卡的布局结构。

系统主要分为三个功能选项卡：
\begin{itemize}
    \item \textbf{智能增强} 负责基础图像文件的导入导出，并提供 CLAHE 增强、伪彩色、Canny 边缘以及 2D-FFT 频谱分析。
    \item \textbf{交互式智能分割：} 提供鼠标手绘框选区域（ROI）的功能，右侧小窗口会实时动态显示框选区域的轮廓提取结果。
    \item \textbf{3D特征映射可视化：} 将二维图像转换为三维空间地貌图，支持用鼠标在画板上进行任意视角的旋转和缩放。
\end{itemize}

% 截图位置 1：软件整体界面布局
\begin{figure}[htbp]
    \centering
    \includegraphics[width=0.85\textwidth]{ui_main.png}
    \caption{系统主界面布局示意图}
    \label{fig:ui_main}
\end{figure}

\section{核心功能与关键代码实现}

\subsection{全局变量与数据管理}
在 App 后台，我定义了私有属性来存储原始图像和当前处理的图像矩阵。这种设计避免了使用全局变量引起的命名冲突，保证了数据在不同选项卡和按钮回调之间的安全传递。

\begin{lstlisting}
properties (Access = private)
    ImgOriginal % 存储读入的原始图像数据
    ImgCurrent  % 存储当前算法处理后的图像数据
end
\end{lstlisting}

\subsection{算法动态参数调节}
为了实现人机交互的实时性，我将滑块与算法参数进行了绑定。当用户拖动滑块时，程序会读取数值并动态计算。以下是滑块数值改变时的核心控制代码：

\begin{lstlisting}
% 滑块数值改变回调函数
function Slider_IntensityValueChanged(app, event)
    if isempty(app.ImgOriginal), return; end
    sliderVal = app.Slider_Intensity.Value; % 获取0-100滑块值
    selectedAlgo = app.Drop_Algo.Value;     % 获取当前下拉菜单算法
    
    switch selectedAlgo
        case 'CLAHE限制对比度增强'
            % 将滑块值动态映射为自适应直方图均衡化的裁剪阈值
            clipLimitValue = 0.01 + (sliderVal / 100) * 0.19; 
            if size(app.ImgOriginal, 3) == 3
                labImg = rgb2lab(app.ImgOriginal);
                L = labImg(:,:,1) / 100;
                L_enhanced = adapthisteq(L, 'NumTiles', [4 4], 'ClipLimit', clipLimitValue);
                labImg(:,:,1) = L_enhanced * 100;
                app.ImgCurrent = lab2rgb(labImg);
            else
                app.ImgCurrent = adapthisteq(app.ImgOriginal, 'NumTiles', [4 4], 'ClipLimit', clipLimitValue);
            end
            
        case '边缘提取'
            % 滑块越大，Canny算子检测阈值越低，提取的边缘细节越丰富
            thresh = 0.1 + ((100 - sliderVal) / 100) * 0.7; 
            if size(app.ImgOriginal, 3) == 3
                grayImg = rgb2gray(app.ImgOriginal);
            else
                grayImg = app.ImgOriginal;
            end
            edgeImg = edge(grayImg, 'Canny', thresh);
            app.ImgCurrent = uint8(edgeImg) * 255;
    end
    imshow(app.ImgCurrent, 'Parent', app.UIAxes_Result); % 刷新右侧画板
end
\end{lstlisting}


% 截图位置 2：算法处理效果对比
\begin{figure}[htbp]
    \centering
    \includegraphics[width=0.85\textwidth]{algorithm_result.png}}
    \caption{图像处理算法前后对比效果}
    \label{fig:algorithm}
\end{figure}

\subsection{实时截图}
在第二个选项卡中，为了实现“鼠标拖拽红框，右侧小图跟着实时刷新”的动态效果，我没有采用静态单次截图的做法，而是引入了事件监听机制，将红框的拖动事件与后台裁剪函数进行绑定。

\begin{lstlisting}
% 开启交互框选按钮回调
function Btn_ROIPushed(app, event)
    if isempty(app.ImgOriginal), return; end
    
    % 锁定画板，防止画框时背景图片被鼠标拖走
    disableDefaultInteractivity(app.UIAxes_SegMain);
    imshow(app.ImgOriginal, 'Parent', app.UIAxes_SegMain);
    
    try
        % 激活鼠标画矩形框工具
        roi = drawrectangle(app.UIAxes_SegMain, 'Label', '目标提取区', 'Color', 'r', 'LineWidth', 2);
        
        % 第一次绘制时手动触发一次小图更新
        app.processAndShowCrop(roi.Position);
        
        % 绑定监听器：只要检测到红框在移动（MovingROI），就自动调用后台刷新函数
        addlistener(roi, 'MovingROI', @(src, event) app.processAndShowCrop(event.CurrentPosition));
    catch 
        disp('用户取消了操作');
    end
end
\end{lstlisting}

% 截图位置 3：手绘ROI实时联动效果
\begin{figure}[htbp]
    \centering
    \includegraphics[width=0.85\textwidth]{roi_live.png}
    \caption{手绘ROI实时联动与边缘提取界面}
    \label{fig:roi}
\end{figure}

\subsection{二维图像的3D亮度地貌重构}
为了更深层次地展示图像的灰度特征，第三个选项卡将二维矩阵转换为三维空间网络。代码通过降低分辨率（防止卡顿）并调用 `surf` 函数，绘制出图像灰度值的拓扑波动图。

\begin{lstlisting}
% 生成3D图按钮回调
function Btn_3DPushed(app, event)
    if isempty(app.ImgOriginal), return; end
    
    % 缩放图像至200x200，保证3D旋转时的流畅度
    if size(app.ImgOriginal, 3) == 3
        grayImg = rgb2gray(app.ImgOriginal);
    else
        grayImg = app.ImgOriginal;
    end
    smallImg = imresize(grayImg, [200, 200]); 
    
    % 绘制三维表面图
    surf(app.UIAxes_3D, double(smallImg), 'EdgeColor', 'none');
    colormap(app.UIAxes_3D, 'parula');
    view(app.UIAxes_3D, -37.5, 30); % 设置初始倾斜视角
    grid(app.UIAxes_3D, 'on');
    enableDefaultInteractivity(app.UIAxes_3D); % 开启鼠标拖动旋转
end
\end{lstlisting}

% 截图位置 4：3D地貌图旋转展示
\begin{figure}[htbp]
    \centering
    \includegraphics[width=0.8\textwidth]{view_3d.png}
    \caption{图像灰度空间3D拓扑结构展示}
    \label{fig:3d}
\end{figure}

\section{调试过程与遇到的问题及解决办法}
在开发这个 App 的过程中，我遇到了几个具体的硬件交互和编译逻辑问题，通过查阅文档和调整代码结构，最终逐一解决：

\begin{enumerate}
    \item \textbf{手绘框选时图片移位问题：}
    在编写第二个选项卡时，我发现用鼠标拖动红框时，背景图片经常会跟着一起漂移，导致裁剪错位。经过排查，发现这是由于 MATLAB 画板默认开启了平移（Pan）手势。
    \\ \textbf{解决办法：} 在画框前加入 \texttt{disableDefaultInteractivity(app.UIAxes\_SegMain)} 语句，锁死背景图，彻底解决了画框和拖图的手势冲突。
    
    \item \textbf{双击空白处或关闭时报错问题：}
    在测试 ROI 裁剪时，如果鼠标不小心在画板上双击或者直接关闭窗口，后台命令行会弹出一大堆关于坐标未定义的红色报错。
    \\ \textbf{解决办法：} 我将整个框选代码包裹在 \texttt{try-catch} 块中。这样即使操作中途被强行中断，程序也会优雅退出，不会报错崩溃。
    
    \item \textbf{添加私有函数时按钮变灰问题：}
    在编写后台实时截图函数时，我发现 Code View 顶部的“添加函数”按钮一直是灰色的，无法点击。
    \\ \textbf{解决办法：} 查阅手册后发现是因为鼠标光标停留位置不对导致系统丢失了焦点。我直接通过手动编写 \texttt{methods (Access = private)} 声明代码块，绕过了图形界面的按钮限制，直接把方法写进了底层的类定义中。
\end{enumerate}

% 截图位置 4：3D地貌图旋转展示
\begin{figure}[htbp]
    \centering
    \includegraphics[width=0.8\textwidth]{alert1.png}
    \caption{系统异常防御与提示弹窗}
    \label{fig:alert1}
\end{figure}

\section{总结与设计收获}
这次基于 MATLAB App Designer 的课程设计，不仅让我巩固了课本上关于数字信号处理和矩阵映射的算法知识，更重要的是锻炼了我的软件工程实践能力。

在设计过程中，我理解了如何通过“私有属性”管理数据流，如何利用“事件监听器”设计出响应快速的交互界面，以及如何通过容错设计来保证软件的稳定性。最后，我成功将项目打包发布为 `.mlappinstall` 格式的安装包，完成了从需求分析、界面排版、代码编写到打包发布的完整流程。这为我后续深入学习电子信息专业的相关软件设计打下了良好的基础。

\end{document}